% DUBTE: les distàncies (10^10 km) no són les reals de Mercuri (~10^7 km); es transcriuen tal com són a l'original.
\begin{apartat}{2,5}
L'òrbita de Mercuri al voltant del Sol és el·líptica i presenta la major excentricitat entre els planetes del sistema solar. Això fa que la seva distància al Sol variï considerablement, afectant la seva velocitat orbital. A l'afeli el planeta Mercuri és a $6,99\cdot 10^{10}$~km del Sol, i al periheli és a $4,60\cdot 10^{10}$~km del Sol. La seva velocitat orbital és de $3,88\cdot 10^{4}$~m/s a l'afeli. Dibuixeu aproximadament l'òrbita de Mercuri i indiqueu la posició de l'afeli i el periheli. Utilitzant el principi de conservació del moment angular, calculeu quin és el mòdul de la velocitat de Mercuri al periheli. Argumenteu si el mòdul de la velocitat varia quan Mercuri es desplaça al llarg de la seva òrbita, i justifiqueu en quin punt de l'òrbita pren el seu valor mínim.
\end{apartat}
